% Preámbulo común de los informes de laboratorio (II/2026). Se incluye desde laboratorios/labNN_*/
%   \input{../informe_comun}
\input{../../gemelo-digital/manual/estilo_manual}
\input{../../gemelo-digital/manual/tecnico/macros}
\graphicspath{{../../documentacion/planos_isa/pid/}{../../documentacion/planos_isa/electrico/}{../../documentacion/manuales_inversa/}
  {../../database/mockups/}{../../database/documentacion/}{../../database/modelo/}{../../database/interfaz/}
  {../../gemelo-digital/manual/img/}{../../modelo-3d/informe/figuras/}{../../modelo-3d/levantamiento/}{../../../../global/fotos-planta/}{../figuras/}}
\usepackage{pdflscape}
\lstset{showstringspaces=false}
% plano A3 en una página apaisada: \plano[página]{archivo}{leyenda}
\newcommand{\plano}[3][1]{\begin{landscape}\begin{figure}[p]\centering\includegraphics[page=#1, width=\linewidth, height=0.86\textwidth, keepaspectratio]{#2}\caption{#3}\end{figure}\end{landscape}}
\setlength{\headheight}{14pt}
\newcommand{\nlab}{}
\newcommand{\tlab}{}
\lhead{\small Brazo XK-JS1-1 · Informe del laboratorio \nlab}
\rhead{\small UMSA · II/2026}

% \portada{número}{título}{inicio}{entrega}{semanas}
\newcommand{\portada}[5]{%
\renewcommand{\nlab}{#1}%
\begin{titlepage}
\centering
\vspace*{0.6cm}
{\large Universidad Mayor de San Andrés · Facultad de Ingeniería\par}
{\large Carrera de Ingeniería Electrónica\par}
\vspace{0.8cm}
{\color{naranja!80!black}\rule{\linewidth}{1.5pt}\par}
\vspace{0.3cm}
{\Large Informe del Laboratorio Integrado Nº #1\par}
\vspace{0.2cm}
{\huge\bfseries #2\par}
\vspace{0.3cm}
{\color{naranja!80!black}\rule{\linewidth}{1.5pt}\par}
\vspace{0.8cm}
{\large Planta: brazo robótico XK-JS1-1 de 4 grados de libertad\par}
\vspace{0.3cm}
{\large ETN-902 Sistemas de Control II · ETN-1034 Aplicación de Técnicas de Control · ETN-1000 Bases de Datos\par}
\vspace{1cm}
\begin{tabular}{rl}
\textbf{Docente:} & Ing. Jorge Antonio Nava Amador\\
\textbf{Auxiliares:} & Limber Ajnota · Fabián Plata · Franz Choque\\
\textbf{Inicio de trabajo:} & #3\\
\textbf{Entrega:} & #4 (#5)\\
\end{tabular}
\vspace{1cm}

\begin{tabular}{ll}
\multicolumn{2}{l}{\textbf{Escuadra del brazo de 4 GDL}}\\[2pt]
Control (ETN-902 / ETN-1034): & Gabriel Andres Gomez Calla\\
 & Jose Lenin Mamani Choque\\
 & Miguel Alejandro Lima Huarachi\\
 & Román Abner Machicado Arequipa\\[2pt]
Cómputo (ETN-1000): & Miguel Alejandro Lima Huarachi\\
 & Marco Eduardo Chura Rondo\\
\end{tabular}
\vfill
{\small Repositorio: \url{https://github.com/GabAndGom/Brazo-de-4-grados-de-libertad}\par}
{\small La Paz, Bolivia · gestión II/2026\par}
\end{titlepage}
}

% Tabla de entregables (Definition of Done): \begin{dod} ... \end{dod}, filas: entregable & estado & evidencia\\
\newenvironment{dod}{\begin{center}\small\begin{longtable}{>{\raggedright}p{5.2cm} >{\raggedright}p{1.8cm} >{\raggedright\arraybackslash}p{7.6cm}}
\toprule \textbf{Entregable} & \textbf{Estado} & \textbf{Evidencia / observación}\\ \midrule \endhead}
{\bottomrule\end{longtable}\end{center}}
